
\documentclass[10pt,twocolumn]{article}

% ================================================================
% IDIOMA, CODIFICACION Y FUENTES
% ================================================================
\usepackage[utf8]{inputenc}
\usepackage[T1]{fontenc}
\usepackage[spanish,es-noshorthands]{babel}
\addto\captionsspanish{\renewcommand{\tablename}{Tabla}}
\addto\captionsspanish{\renewcommand{\refname}{Referencias}}

\usepackage{amsmath}
\usepackage{newtxtext,newtxmath}
\usepackage{microtype}

% ================================================================
% PAGINA Y FORMATO
% ================================================================
\usepackage[letterpaper,margin=2.2cm,columnsep=0.8cm]{geometry}
\usepackage{titlesec}
\usepackage{booktabs,tabularx,array}
\usepackage{enumitem}
\usepackage{xurl}
\usepackage[hidelinks]{hyperref}

\setlist[itemize]{leftmargin=1.2em,itemsep=2pt,topsep=3pt}
\setlength{\parindent}{1.2em}
\setlength{\parskip}{0pt}
\emergencystretch=2em
\raggedbottom

\titleformat{\section}{\Large\bfseries}{\thesection}{0.6em}{}
\titleformat{\subsection}{\large\bfseries\itshape}{\thesubsection}{0.6em}{}
\titleformat{\subsubsection}{\normalsize\bfseries}{\thesubsubsection}{0.6em}{}

\titlespacing*{\section}{0pt}{1.6ex plus .5ex}{1ex}
\titlespacing*{\subsection}{0pt}{1.3ex plus .4ex}{0.8ex}
\titlespacing*{\subsubsection}{0pt}{1.1ex plus .3ex}{0.6ex}

\begin{document}

% ================================================================
% ENCABEZADO
% ================================================================
\twocolumn[
\begin{@twocolumnfalse}
\raggedright

{\LARGE\bfseries
Diseño de un sistema de anonimización de voz en tiempo real para proteger
la identidad vocal en telefonía móvil en Colombia\par}

\vspace{10pt}
{\large Alison Daiana Aristizabal Garcia$^{1,2}$\par}
\vspace{2pt}
{\large Asesores: Juan Fernando Pérez$^{1}$ y Yezid Donoso$^{2}$\par}
\vspace{8pt}
{\small
$^{1}$ Departamento de Ingeniería Industrial, Universidad de los Andes\par
$^{2}$ Departamento de Ingeniería de Sistemas y Computación, Universidad de los Andes\par}

\vspace{8pt}
\rule{\textwidth}{0.4pt}
\vspace{6pt}

\noindent\textbf{Abstract}\par

\vspace{6pt}
\noindent\textbf{Index Terms:}

\vspace{10pt}
\noindent\textbf{Resumen}\par

\vspace{6pt}
\noindent\textbf{Palabras clave:}

\vspace{6pt}
\rule{\textwidth}{0.4pt}
\vspace{8pt}

\end{@twocolumnfalse}
]

% ================================================================
% 1. INTRODUCCION
% ================================================================
\section{Introducción}

El fraude telefónico y el uso de técnicas de inteligencia artificial (IA) para suplantar la identidad de las personas son un problema creciente de seguridad y privacidad. La voz, además de ser uno de los principales medios de comunicación, contiene características biométricas que permiten distinguir a un hablante. Esas características pueden ser aprovechadas tanto por sistemas de reconocimiento como por tecnologías de síntesis y clonación de voz. Barrington et al.~\cite{barrington2025} mostraron que las personas no identifican de manera fiable las voces generadas con IA: en su estudio, los participantes juzgaron que una voz clonada pertenecía a la misma persona que la voz original alrededor del 80\,\% de las veces, y solo reconocieron como artificial cerca del 60\,\% de las voces clonadas. Informes recientes sobre fraude describen además el uso de la clonación de voz en esquemas de ingeniería social \cite{groupib2025}.

En Colombia, el problema ocurre en un contexto de alta exposición a llamadas no deseadas. Según datos de la aplicación Truecaller reseñados por ACIS, durante 2025 los colombianos recibieron más de 16\,630 millones de llamadas spam \cite{acis2026}. En una encuesta de TransUnion Colombia a 995 personas, el 27\,\% de quienes afirmaron haber sido blanco de un intento de fraude digital entre agosto y diciembre de 2024 indicó que el intento ocurrió mediante \textit{vishing} (\textit{voice phishing}), es decir, una llamada en la que el estafador busca que la víctima entregue información personal; fue la modalidad más reportada \cite{transunion2025}. En este escenario, cuando una persona contesta una llamada de un número desconocido, fragmentos de su voz quedan expuestos a un interlocutor potencialmente malicioso.

Las soluciones disponibles se concentran en identificar, clasificar o bloquear números sospechosos. Estos mecanismos reducen el contacto con llamadas no deseadas, pero no protegen la voz una vez que la conversación ha comenzado. Este vacío motiva el estudio de mecanismos preventivos que actúen sobre la voz del usuario durante la comunicación.

El presente proyecto diseña y evalúa un sistema de anonimización de voz en tiempo real. Anonimizar la voz significa transformarla para que sea más difícil reconocer quién habla, sin cambiar lo que se dice. La propuesta se basa en una transformación estocástica mediante el coeficiente de McAdams \cite{patino2021}, seleccionada por tres razones: su bajo costo computacional, la posibilidad de ejecutarla en la unidad central de procesamiento (CPU, \textit{Central Processing Unit}) de un computador personal sin una tarjeta gráfica, y la ausencia de una etapa de entrenamiento con grandes corpus de voz.

El sistema tiene dos componentes. El primero es un \textit{pipeline} por lotes, que anonimiza un corpus completo de grabaciones y calcula las métricas de evaluación. El segundo es un prototipo de escritorio en tiempo real, que anonimiza la voz capturada por el micrófono mientras la persona habla. El alcance del proyecto llega hasta la anonimización de la voz en tiempo real.

Para evaluar el sistema se usan tres métricas, cada una asociada con una parte del problema. La privacidad se mide con la tasa de error (EER, \textit{Equal Error Rate}) de un sistema de verificación automática de locutor (ASV, \textit{Automatic Speaker Verification}), que actúa como un atacante que intenta reconocer al hablante. La inteligibilidad se mide con la tasa de error de palabras (WER, \textit{Word Error Rate}) de un sistema de reconocimiento automático de habla (ASR, \textit{Automatic Speech Recognition}), que comprueba si el contenido se sigue entendiendo. El tiempo real se mide con el factor de tiempo real (RTF, \textit{Real-Time Factor}), que indica si el procesamiento es lo bastante rápido para una conversación. Estos conceptos se definen en detalle en la Sección~2.

La pregunta de investigación que orienta el trabajo es: ¿en qué medida un sistema de anonimización de voz en tiempo real, activado ante llamadas de números desconocidos, protege la identidad vocal de usuarios de telefonía móvil en Colombia sin comprometer la inteligibilidad de la comunicación?

\subsection{Justificación}

El proyecto responde a un riesgo asociado con la exposición de la identidad vocal durante llamadas telefónicas. La elevada cantidad de llamadas spam y la participación de las llamadas en los intentos de fraude digital en Colombia muestran un escenario de exposición relevante para los usuarios \cite{acis2026,transunion2025}. Al mismo tiempo, la disponibilidad de tecnologías capaces de generar voces artificiales convincentes aumenta la importancia de proteger las características biométricas de la voz \cite{barrington2025,groupib2025}.

El valor agregado del proyecto consiste en abordar el problema en el momento en que ocurre la exposición de la voz. En lugar de limitarse a identificar números potencialmente fraudulentos, se estudia la viabilidad de transformar la voz antes de que llegue al interlocutor, reduciendo la información asociada con la identidad del hablante y conservando el contenido de la conversación. Esto plantea un problema de ingeniería en el que deben equilibrarse a la vez la privacidad, la inteligibilidad y el desempeño computacional.

Además, el proyecto aporta evidencia experimental sobre anonimización de voz en español. Los trabajos comparados en la Sección~3 se evaluaron en inglés o en escenarios distintos al de este proyecto. Este trabajo aporta evidencia sobre el comportamiento de una técnica de procesamiento de señal de bajo costo computacional en un contexto de interés para usuarios de telefonía móvil en Colombia.

% ================================================================
% 2. MARCO TEORICO
% ================================================================
\section{Marco teórico}

\subsection{Voz, identidad vocal y datos biométricos}

La señal de voz transporta el contenido lingüístico, es decir, qué se dice, y rasgos asociados con cómo se dice. Entre ellos está la huella vocal (\textit{voiceprint}), que permite distinguir a un hablante de otro \cite{deng2023}. Por esta razón, la literatura de privacidad de voz considera la voz una fuente de información biométrica que permite identificar a una persona \cite{patino2021}.

Las características biométricas de la voz se relacionan con propiedades fisiológicas y de comportamiento del hablante, como la frecuencia fundamental, los formantes, el timbre, la prosodia y el ritmo del habla. La combinación de estos rasgos permite que un sistema automático compare dos grabaciones y estime si pertenecen a la misma persona.

Por esta razón, proteger la voz durante una llamada exige reducir esa información sin eliminar el mensaje, que es la tarea de la anonimización descrita en la Sección~2.2.

\subsection{Anonimización y seudonimización de voz}

Anonimizar voz consiste en modificar una señal de habla para reducir la capacidad de identificar al hablante, conservando propiedades relevantes de la comunicación, como la inteligibilidad y la naturalidad \cite{patino2021}. El \textit{VoicePrivacy Challenge} (VPC) es una iniciativa internacional que define cómo evaluar estos sistemas: el sistema debe producir una nueva forma de onda, reducir la información del hablante y preservar la utilidad de la señal \cite{tomashenko2024}.

Como la voz transformada puede seguir siendo distinguible de otras voces sin corresponder a la identidad original, Patino et al.~\cite{patino2021} señalan que la tarea se describe con más precisión como seudonimización. En este trabajo se usa el término anonimización por consistencia con la literatura técnica del área.

El efecto de una anonimización no se supone: se mide con la EER y la WER (Secciones~2.3 y 2.4), siguiendo los escenarios de ataque del VPC (Sección~2.5). Una anonimización exitosa en este sentido no garantiza que la identidad sea imposible de recuperar.

\subsection{Verificación automática de locutor (ASV) y EER}

Un sistema de verificación automática de locutor (ASV) decide si dos grabaciones pertenecen a la misma persona. El proceso tiene dos etapas. En la inscripción (\textit{enrollment}), el sistema recibe grabaciones de un hablante y construye una representación de su voz. En la prueba (\textit{trial}), recibe una grabación nueva, la compara con esa representación y obtiene una puntuación de similitud \cite{deng2023}. Si la puntuación supera un umbral, el sistema acepta que es la misma persona.

Con cualquier umbral, el sistema puede cometer dos errores. La tasa de falsa aceptación (FAR, \textit{False Acceptance Rate}) es la proporción de comparaciones entre personas distintas que el sistema acepta como la misma persona. La tasa de falso rechazo (FRR, \textit{False Rejection Rate}) es la proporción de comparaciones de una misma persona que el sistema rechaza. La EER es el valor en el que ambas tasas son iguales:

\begin{equation}
    \mathrm{FAR}(\tau) = \mathrm{FRR}(\tau),
    \label{eq:eer}
\end{equation}

donde $\tau$ es el umbral de decisión. Una EER cercana a 0\,\% indica que el sistema reconoce bien a los hablantes; una EER cercana a 50\,\% indica que acierta como si decidiera al azar.

Cuando el ASV recibe grabaciones anonimizadas, cumple el papel de un atacante que intenta reconocer al hablante. Si la EER es alta, ese atacante se equivoca con frecuencia; por eso, al evaluar una anonimización, una EER más alta significa más privacidad. Para este fin se emplea el modelo preentrenado ECAPA-TDNN (\textit{Emphasized Channel Attention, Propagation and Aggregation in Time Delay Neural Network}) \cite{desplanques2020}, sin modificarlo.

\subsection{Reconocimiento automático de habla (ASR) y WER}

Un sistema de reconocimiento automático de habla (ASR) transcribe el contenido de una grabación, sin importar quién habla. La inteligibilidad se evalúa con la tasa de error de palabras (WER):

\begin{equation}
    \mathrm{WER}
    =
    \frac{N_{\mathrm{sust}} + N_{\mathrm{elim}} + N_{\mathrm{ins}}}
         {N_{\mathrm{ref}}},
    \label{eq:wer}
\end{equation}

donde $N_{\mathrm{sust}}$, $N_{\mathrm{elim}}$ y $N_{\mathrm{ins}}$ son el número de palabras sustituidas, eliminadas e insertadas respecto a la transcripción de referencia, y $N_{\mathrm{ref}}$ es el número de palabras de esa referencia \cite{deng2023}. Una WER más baja indica que el contenido se conserva mejor. Como las inserciones también cuentan como errores, la WER puede superar el 100\,\% cuando el ASR agrega muchas palabras que no se dijeron.

Al comparar la WER de las grabaciones originales con la de las anonimizadas se puede saber si el mensaje se sigue entendiendo: si la WER aumenta poco, la anonimización conservó el contenido. Para transcribir se emplea el modelo preentrenado Whisper, versión \textit{turbo}, basado en el trabajo de Radford et al.~\cite{radford2023}, sin modificarlo.

\subsection{Modelos de atacante}

El resultado de una evaluación de anonimización depende de lo que se supone que sabe el atacante. El VPC contempla varios escenarios \cite{patino2021,tomashenko2024}. Cada uno se nombra con dos letras: la primera indica cómo están las grabaciones de inscripción y la segunda cómo están las de prueba, con ``o'' para original y ``a'' para anonimizada.

\begin{itemize}
    \item \textbf{Atacante ignorante (o--a):} no sabe que el audio fue anonimizado. Inscribe al hablante con grabaciones originales y las compara con grabaciones anonimizadas.

    \item \textbf{Atacante semi-informado (a--a):} conoce el sistema de anonimización, pero no la configuración exacta usada en las grabaciones de prueba. Anonimiza también sus grabaciones de inscripción y compara audio anonimizado con audio anonimizado (escenario anonimizado--anonimizado).

    \item \textbf{Atacante informado:} conoce el algoritmo y su configuración. Patino et al.~\cite{patino2021} no incluyen este escenario en su evaluación principal.
\end{itemize}

Quamer y Gutierrez-Osuna~\cite{quamer2024} usan además el término \textit{lazy-informed} para el caso en que la inscripción y la prueba se anonimizan con hablantes objetivo distintos.

A estos escenarios se suma el o--o (original--original), sin anonimización, que sirve de referencia: indica qué tan bien reconoce el ASV a los hablantes sin protección. En el contexto de una llamada, el escenario o--a representa a un atacante que graba la conversación y la compara con grabaciones reales de la víctima, y el a--a a un atacante que conoce el sistema y lo usa para anonimizar sus propias grabaciones de la víctima.

\subsection{Modelo fuente--filtro y codificación predictiva lineal (LPC)}

El modelo fuente--filtro describe cómo se produce la voz. Se puede comparar con una trompeta: el músico hace vibrar los labios contra la boquilla y produce un zumbido, y el tubo de metal transforma ese zumbido en el sonido de la trompeta. En la voz, la fuente es el sonido que se produce en la garganta y el filtro es el tracto vocal, formado por la garganta, la boca y la nariz. Este modelo separa la voz en dos componentes: la excitación, que corresponde a la fuente, y los formantes, que dependen de la forma del filtro.

\subsubsection{Excitación}

La excitación es el sonido que produce la fuente antes de pasar por el tracto vocal. Proviene de la vibración de las cuerdas vocales. La excitación lleva la frecuencia fundamental, es decir, qué tan grave o aguda suena la voz, y el ritmo y la entonación del habla.

Si la excitación se conserva sin cambios, se conservan también el ritmo y la entonación con los que se dice el mensaje, lo que ayuda a que siga siendo comprensible. Esto implica también que la frecuencia fundamental del hablante no cambia.

\subsubsection{Formantes}

Al pasar por el tracto vocal, la excitación se transforma: el tubo refuerza algunas frecuencias y atenúa otras. Las zonas de frecuencia reforzadas se llaman formantes. Su posición depende de la forma y el tamaño del tracto vocal. Por eso, los formantes cambian con cada sonido que se pronuncia, ya que la boca toma una forma distinta en cada vocal, y también varían de una persona a otra, porque cada persona tiene un tracto vocal diferente. Esta segunda propiedad hace que los formantes contribuyan a que cada voz suene distinta.

Por esta razón, desplazar los formantes hace que la voz deje de sonar como el tracto vocal del hablante original.

\subsubsection{Codificación predictiva lineal (LPC)}

Para modificar los formantes sin tocar la excitación, primero hay que separarlos. La codificación predictiva lineal (LPC, \textit{Linear Predictive Coding}) es la técnica matemática que hace esa separación a partir de una grabación.

La idea de la LPC es que cada muestra de una señal de voz se puede predecir bastante bien como una combinación de las $p$ muestras anteriores:

\begin{equation*}
    \hat{s}[n]
    = a_1 s[n-1] + a_2 s[n-2] + \cdots + a_p s[n-p].
\end{equation*}

Los pesos $a_1,\ldots,a_p$ se llaman coeficientes LPC y se eligen para que el error de predicción sea el menor posible. Esos coeficientes describen el filtro, es decir, la forma del tracto vocal. Lo que la predicción no logra explicar, el error $s[n]-\hat{s}[n]$, se llama residuo de predicción y corresponde a la excitación.

En términos técnicos, la LPC representa el tracto vocal como un filtro todo-polo de orden $p$. Los polos son las raíces del polinomio que forman los coeficientes. Los polos complejos aparecen en pares, y cada par cercano al círculo unitario corresponde a un formante: el ángulo del polo indica la frecuencia del formante y el radio indica qué tan marcado es. Por esta relación, mover el ángulo de los polos equivale a mover los formantes.

Como la forma de la boca cambia todo el tiempo al hablar, la LPC no se calcula una sola vez sobre toda la grabación, sino sobre fragmentos cortos de audio, llamados tramas, en los que el tracto vocal casi no cambia. El orden $p$ fija cuántos coeficientes describen el filtro. Una heurística comúnmente atribuida a Markel y Gray \cite{markel1976linear} define el orden como:

\begin{equation}
    p = f_s + 4,
    \label{eq:lpc-order}
\end{equation}

donde $f_s$ es la frecuencia de muestreo expresada en kHz. Para una señal de 16~kHz, esta heurística produce $p=20$.

El orden $p$ es el número de coeficientes, y por tanto de polos, del filtro. Como cada formante necesita un par de polos y en el habla aparece aproximadamente un formante por cada kHz del rango de frecuencias de la señal (hasta $f_s/2$), se requieren $f_s$ polos para los formantes, más 4 para la forma general del espectro de la fuente: a 16~kHz, $p=20$. Elegir bien $p$ es importante. Si es muy pequeño, algunos formantes quedan sin representar y se mezclan con la excitación; si es muy grande, el filtro describe detalles que no son formantes, como los armónicos de la voz. En ambos casos, la separación entre formantes y excitación deja de ser correcta.

Al calcular la LPC en cada trama, con el orden $p$ de la Ecuación~\eqref{eq:lpc-order}, se obtienen por separado el filtro, representado por sus polos, y la excitación. Si se modifica solo el ángulo de los polos y luego se pasa la excitación original por el filtro modificado, los formantes se desplazan: cambian los rasgos del tracto vocal del hablante y se conserva lo que se dice. La Sección~2.7 explica cómo se modifican los polos.

\subsection{Coeficiente de McAdams}

Patino et al.~\cite{patino2021} usan el coeficiente de McAdams para transformar la envolvente espectral modificando la posición angular de los polos LPC. El efecto principal es un desplazamiento de los formantes, sin necesidad de entrenar un modelo.

Un polo complejo se expresa como:

\begin{equation}
    z_i = r_i\,e^{j\phi_i},
    \label{eq:polo}
\end{equation}

donde $r_i$ es el radio y $\phi_i$ la posición angular, medida en radianes entre $0$ y $\pi$. El ángulo indica la frecuencia del formante asociado al polo: $f_i=\phi_i f_s/(2\pi)$.

La transformación de McAdams eleva el ángulo de cada polo complejo a la potencia $\alpha$ y deja su radio sin cambios:

\begin{equation}
    z_i' = r_i\,e^{j\phi_i^{\alpha}},
    \label{eq:mcadams-transform}
\end{equation}


\begin{equation}
    \alpha \sim \mathcal{U}(\alpha_{\min},\alpha_{\max}),
    \label{eq:alpha}
\end{equation}

donde $\mathcal{U}(a,b)$ indica que cualquier valor entre $a$ y $b$ tiene la misma probabilidad de ser elegido. Cambiar el intervalo de $\alpha$ modifica el compromiso entre privacidad e inteligibilidad \cite{patino2021}.

En el sistema propuesto, $\alpha$ se sortea una vez por llamada y se mantiene fijo durante toda la llamada; en la evaluación sobre un corpus de grabaciones, se sortea un $\alpha$ por grabación. Como el intervalo determina el equilibrio entre privacidad e inteligibilidad, se trata como un parámetro de diseño que se evalúa experimentalmente, no como un valor fijado de antemano.

\subsubsection{ZEBRA: evaluación de evidencia biométrica}

ZEBRA (\textit{Zero Evidence Biometric Recognition Assessment}) es un marco de evaluación que mide cuánta evidencia sobre la identidad del hablante permanece en las grabaciones anonimizadas, y la resume en categorías. La categoría ``0'' corresponde a evidencia nula, es decir, a privacidad perfecta. Patino et al.~\cite{patino2021} reportan que ninguna configuración basada en el coeficiente de McAdams alcanza esa categoría.

Aunque ZEBRA no se calcula en este trabajo, el resultado de Patino et al.~\cite{patino2021} deja claro que el coeficiente de McAdams reduce la información del hablante, pero no logra una anonimización perfecta.

\subsection{Tiempo real: latencia y factor de tiempo real (RTF)}

El factor de tiempo real (RTF) es la razón entre el tiempo que tarda el sistema en procesar un audio y la duración de ese audio:

\begin{equation}
    \mathrm{RTF}
    =
    \frac{T_{\text{procesamiento}}}{T_{\text{audio}}}.
    \label{eq:rtf}
\end{equation}

Un RTF menor que 1 indica que el sistema procesa más rápido de lo que dura el audio. Por ejemplo, un RTF de 0.1 significa que procesar 10~s de voz toma 1~s. Deng et al.~\cite{deng2023} usan la misma definición con el nombre de coeficiente de tiempo real (RTC, \textit{Real-Time Coefficient}). La latencia es el retardo entre el momento en que se habla y el momento en que sale la voz transformada; Quamer y Gutierrez-Osuna~\cite{quamer2024} la calculan como el tamaño del fragmento de audio más el tiempo de cómputo.

Por eso, el RTF del núcleo de anonimización permite comprobar si el procesamiento es lo bastante rápido para una conversación. La latencia algorítmica, por su parte, depende del tamaño de los fragmentos de audio que se procesan (Sección~4.4).

% ================================================================
% 3. METODOLOGIA
% ================================================================
\section{Metodología}

La metodología comenzó con la formulación y delimitación del problema. Luego se revisó la viabilidad 
de implementar el proyecto y se definieron las restricciones de diseño. Con esas restricciones como 
referencia, se compararon métodos alternativos de anonimización y se seleccionó el que mejor las cumplía. 
La implementación empezó por el núcleo de procesamiento y continuó con el pipeline por lotes. Después se 
definió e implementó el protocolo de evaluación, que permitió iniciar la experimentación. Finalmente, se 
completó el prototipo de anonimización en tiempo real.

\subsection{Formulación y delimitación del problema}
El problema se formula así: durante una llamada, un tercero puede obtener fragmentos de la voz de una persona 
y usarlos después para identificarla o suplantarla. El sistema propuesto busca que la voz que sale del usuario 
conserve el mensaje, pero que no sirva para reconocer a quien habla.

La pregunta de investigación es: ¿en qué medida un sistema de anonimización de voz en tiempo real, activado 
automáticamente ante llamadas de números desconocidos, protege la identidad vocal de usuarios de telefonía móvil 
en Colombia sin comprometer la inteligibilidad de la comunicación? La activación ante números desconocidos describe 
el escenario de uso que motiva la pregunta; este trabajo evalúa la protección que ofrece la anonimización 
una vez activada.

El objetivo general es diseñar y evaluar un prototipo de sistema de anonimización de voz en tiempo real para 
la protección de la identidad vocal de usuarios de telefonía móvil en Colombia ante llamadas de números 
desconocidos, mediante técnicas de procesamiento de señal y métricas estandarizadas de privacidad e 
inteligibilidad.

\subsection{Viabilidad y restricciones de diseño}
El sistema está pensado para funcionar, en último término, en el teléfono celular del usuario, que es el 
dispositivo donde ocurren las llamadas. Un celular tiene menos capacidad de procesamiento, menos memoria y 
menos energía disponible que un computador, y normalmente no cuenta con una GPU que pueda dedicarse a ejecutar 
modelos de voz. Además, el procesamiento de la voz debe ocurrir durante la llamada, junto con las demás tareas 
del teléfono, y sin agotar la batería. Por eso, la viabilidad del proyecto no se evaluó solo frente al computador 
de desarrollo, sino frente a las limitaciones de un dispositivo móvil.

Con este propósito, el método de anonimización debía cumplir cinco criterios. Cada uno responde a una 
restricción concreta:

\begin{itemize}
    \item Bajo costo computacional y operación en tiempo real en CPU. El sistema debe transformar la voz mientras 
    la persona habla, sin depender de una GPU, y con un costo lo bastante bajo para que pueda ejecutarse en un 
    dispositivo con recursos restringidos como un celular. Por eso el criterio de eficiencia del proyecto 
    (RTF < 0.1) exige que el procesamiento tome como máximo una décima parte de la duración del audio, lo que 
    deja margen para equipos con menos capacidad que el computador de desarrollo.
    \item No requerir entrenamiento con corpus grandes. Entrenar un modelo de voz exige grandes volúmenes de datos 
    y varias GPU durante días, recursos y tiempo que no están al alcance del proyecto. Además, 
    los métodos que dependen de redes neuronales grandes suelen requerir, también al ejecutarse, más memoria 
    y cómputo de los que ofrece un celular.
    \item Ser implementable dentro del tiempo del proyecto de grado, en tiempo y complejidad.
    \item Contar con evidencia publicada bajo las métricas del VPC, para tener una referencia del comportamiento 
    esperado.
    \item Contar con una implementación de referencia pública, para poder verificar la implementación propia.
\end{itemize}

El desarrollo y todas las pruebas del sistema se realizaron sobre la CPU de un computador de escritorio, 
no en un teléfono móvil. Por lo tanto, los resultados muestran que el procesamiento tiene un costo computacional 
bajo en ese equipo y sugieren que habría margen para ejecutarlo en dispositivos de menor capacidad, pero no 
demuestran su funcionamiento en un teléfono. Esa verificación queda como trabajo futuro.

\subsection{Comparación de alternativas}
Se compararon cuatro enfoques de anonimización:
\begin{itemize}
    \item El coeficiente de McAdams estocástico \cite{patino2021}, que desplaza los formantes modificando los polos LPC, sin entrenamiento.
    \item V-Cloak \cite{deng2023}, que añade a la voz perturbaciones adversariales generadas por una red neuronal para engañar a los sistemas ASV.
    \item El modelo de streaming end-to-end de Quamer y Gutierrez-Osuna \cite{quamer2024}, una red neuronal que reemplaza la identidad del hablante con baja latencia.
    \item Stream-Voice-Anon \cite{kuzmin2026}, que combina un códec neuronal de audio con un modelo de lenguaje para generar una voz anonimizada.
\end{itemize}
\begin{table*}[t]
\centering
\caption{Alternativas evaluadas. Las cifras corresponden a los valores reportados por cada trabajo y no son comparables entre sí.}
\label{tab:alternativas}
\small
\renewcommand{\arraystretch}{1.18}
\begin{tabularx}{\textwidth}{
>{\raggedright\arraybackslash}p{3.0cm}
>{\raggedright\arraybackslash}p{2.7cm}
>{\raggedright\arraybackslash}p{3.2cm}
>{\raggedright\arraybackslash}X
>{\raggedright\arraybackslash}p{2.3cm}}
\toprule
\textbf{Enfoque} & \textbf{Entrenamiento / hardware} & \textbf{Tiempo real reportado} & \textbf{Privacidad y utilidad reportadas} & \textbf{Idioma evaluado} \\
\midrule
McAdams estocástico \cite{patino2021}
& Ninguno
& RTC 0.030, medido por Deng et al.~\cite{deng2023} con $\alpha=0.8$
& EER 33.8\,\% en escenario semi-informado para $\mathcal{U}(0.7,0.9)$; WER 13.7\,\% para $\mathcal{U}(0.8,0.9)$
& Inglés \\

Streaming \textit{end-to-end} \cite{quamer2024}
& LibriTTS, 2 GPU V100
& Latencia de 230~ms en versión base y 66~ms en versión \textit{lite} sobre CPU
& EER a--a de 39.2--42.0\,\%; WER de 5.1--6.5\,\%
& Inglés \\

Códec neuronal y modelo de lenguaje \cite{kuzmin2026}
& 8 GPU H100
& 180~ms y RTF 0.93 en GPU portátil; sin tiempo real en CPU
& EER 46.5--47.7\,\% en \textit{lazy-informed} y 18.6--19.0\,\% en semi-informado; WER 4.7--6.6\,\%
& Inglés \\

Ejemplos adversariales \cite{deng2023}
& VoxCeleb1, 2 GPU RTX 3090
& RTC 0.011
& EER 46.10\,\% en escenario ignorante; WER 7.65\,\%
& ASV: inglés y chino. ASR: inglés, chino, francés e italiano \\
\bottomrule
\end{tabularx}
\end{table*}

La Tabla~\ref{tab:alternativas} resume las características y el desempeño
que reporta cada trabajo: entrenamiento y hardware, tiempo real,
privacidad, utilidad e idioma evaluado. Esta información permite
revisar cada enfoque frente a los cinco criterios de la
Sección~3.2. La EER y la WER se presentan según el escenario
de atacante que considera cada trabajo.

Ninguno de los cuatro enfoques fue evaluado en español sobre
audio de telefonía, y ninguno mide resistencia a la clonación
de voz: todas las evaluaciones se hacen frente a sistemas ASV.

Además, las cifras de la Tabla~\ref{tab:alternativas} no son
directamente comparables entre sí, por tres razones principales.

Primero, cambia el modelo de atacante usado en cada evaluación.
La EER se reporta frente a un atacante ignorante en Deng et
al.~\cite{deng2023}, frente a uno semi-informado en Patino et
al.~\cite{patino2021} y Kuzmin et al.~\cite{kuzmin2026}, y en
el escenario anonimizado--anonimizado (a--a) en Quamer y
Gutierrez-Osuna~\cite{quamer2024}. Para un mismo sistema, el
resultado puede variar considerablemente según este supuesto.
Por ejemplo, Kuzmin et al.~\cite{kuzmin2026} reportan una EER
de 46.5--47.7\,\% frente a un atacante \textit{lazy-informed},
que anonimiza las grabaciones de inscripción y de prueba con
hablantes objetivo distintos, pero de 18.6--19.0\,\% frente
a uno semi-informado.

Segundo, cambian los sistemas de verificación de locutor,
los modelos ASR, los datos usados para entrenarlos y los
corpus de evaluación. Patino et al.~\cite{patino2021} muestran
que entrenar los modelos con audio anonimizado puede modificar
de manera importante los resultados. En su configuración
$\mathcal{U}(0.5,0.9)$, el coeficiente $\alpha$ se elige al azar
entre 0.5 y 0.9 siguiendo una distribución uniforme, es decir,
cualquier valor de ese intervalo tiene la misma probabilidad
de ser elegido. Con esa configuración, la WER disminuye de
aproximadamente 45\,\% a 9.6\,\% cuando se usa un sistema ASR
adaptado al audio anonimizado. Asimismo, el método de McAdams
obtiene una WER de 26.54\,\% cuando se evalúa con los sistemas
ASR de Deng et al.~\cite{deng2023}, frente a 13.7\,\% en la
configuración experimental de Patino et al.~\cite{patino2021}.

Tercero, el desempeño temporal se mide con definiciones y
plataformas de hardware diferentes. Deng et al.~\cite{deng2023}
usan el cociente entre el tiempo de procesamiento y la duración
del audio, que denominan RTC; Quamer y
Gutierrez-Osuna~\cite{quamer2024} reportan la latencia como la
suma del tamaño del fragmento y el tiempo de cómputo sobre CPU;
y Kuzmin et al.~\cite{kuzmin2026} informan tanto la latencia
como el RTF sobre una GPU portátil.

Por estas razones, la Tabla~\ref{tab:alternativas} se usa
para ubicar el orden de magnitud de los resultados y los
requerimientos de cada enfoque (entrenamiento, hardware,
idioma y capacidad de operar en tiempo real), y no para
establecer una clasificación numérica directa entre métodos.
La comparación se basa en el cumplimiento de los criterios
de la Sección~3.2 y no en el desempeño numérico. Por la
misma razón, los resultados de este trabajo, obtenidos con
ECAPA-TDNN, Whisper y Common Voice en español, se analizan
frente a los criterios del proyecto y no se comparan
numéricamente con los valores reportados en estos artículos.



\subsection{Selección del método}

El coeficiente de McAdams con $\alpha$ estocástico es el único
de los cuatro enfoques que cumple los cinco criterios de la
Sección~3.2. No requiere entrenamiento con corpus grandes,
se ejecuta en tiempo real sobre CPU, su implementación es de
complejidad moderada y viable para un proyecto de pregrado,
fue evaluado con el protocolo del VPC y cuenta con una
implementación de referencia pública en la línea base B2
del VPC.

Los otros tres enfoques no cumplen al menos uno de los
criterios:

\begin{itemize}
    \item Los tres requieren entrenar redes neuronales, lo
    que incumple el criterio de no requerir entrenamiento
    con corpus grandes.

    \item Stream-Voice-Anon no opera en tiempo real solo
    con CPU, según sus propios autores
    \cite{kuzmin2026}, lo que incumple el criterio de
    operación en tiempo real sobre CPU.

    \item En el caso de Quamer y Gutierrez-Osuna
    \cite{quamer2024}, no se encontró código público,
    lo que incumple el criterio de contar con una
    implementación de referencia.

    \item En el caso de V-Cloak \cite{deng2023}, el
    repositorio no ofrecía el modelo preentrenado al
    momento de la revisión, lo que también incumple
    el criterio de contar con una implementación de
    referencia.
\end{itemize}

El coeficiente $\alpha$ se sortea una vez por sesión,
es decir, por llamada; se mantiene fijo durante toda
la sesión y cambia entre sesiones.

El intervalo de $\alpha$ se trata como un parámetro
de diseño que se elige empíricamente
(Sección~3.6.3). El punto de partida es
$\mathcal{U}(0.7,0.9)$, uno de los intervalos
explorados por Patino et al.~\cite{patino2021}.
Esta notación indica que $\alpha$ se elige al azar
entre 0.7 y 0.9 siguiendo una distribución uniforme,
es decir, cualquier valor de ese intervalo tiene la
misma probabilidad de ser elegido.

\subsection{Estrategia de implementación}

El sistema se construyó por partes, comprobando que cada una
funcionara bien antes de agregar la siguiente.

Por eso, antes de anonimizar cualquier voz, se verificó que
el núcleo de procesamiento no alterara la señal por sí mismo.
Cada etapa se probó primero sin aplicar ninguna transformación:
el audio se dividía en tramas, se analizaba con LPC y se
volvía a reconstruir, y el resultado tenía que ser igual
al audio original. Solo cuando todas las etapas superaron
esta prueba se añadió la transformación de McAdams.
De esta forma, cualquier cambio que aparezca en la voz
anonimizada se debe a la transformación y no a un error
en el resto del sistema.

Después se adaptó el núcleo para procesar el audio por
fragmentos, como ocurre durante una llamada, en la que
la voz llega poco a poco y no está completa desde el inicio.
Para verificarlo, se comparó el resultado de procesar una
grabación completa de una sola vez con el de procesarla
en fragmentos de distintos tamaños. Ambos resultados
debían ser iguales. Por último, se construyó el
\textit{pipeline} por lotes, que aplica el sistema sobre
todas las grabaciones del corpus de evaluación.

Un aspecto importante del diseño es que el \textit{pipeline}
por lotes y el prototipo en tiempo real usan el mismo núcleo
de procesamiento. Así, lo que se mide sobre el corpus
corresponde exactamente a lo que se ejecuta durante una
llamada simulada en vivo, y no a una versión distinta
del sistema.

En cada etapa se escribieron pruebas unitarias, y solo
se pasaba a la siguiente cuando todas eran exitosas.

\subsection{Protocolo de evaluación}
\subsubsection{Datos y muestreo}

La evaluación usa Mozilla Common Voice 26.0 en español
\cite{ardila2020}, y solo los clips del archivo
\texttt{validated.tsv}, es decir, grabaciones que la
comunidad revisó y aceptó. Los términos de uso del corpus
piden no intentar determinar la identidad de los hablantes.
En este trabajo, la reidentificación ocurre solo dentro
del corpus: se comparan grabaciones entre sí para medir
la anonimización, sin vincularlas con identidades reales.

Se aplicaron dos filtros: cada clip debía durar al menos
$3000\,\mathrm{ms}$ y cada hablante debía tener al menos
5 clips. Luego se hizo un muestreo estratificado por acento
autorreportado, con semilla 42, $K=134$ hablantes por
estrato y $N=5$ clips por hablante. El valor de $K$ está
limitado por el estrato más pequeño. La muestra final tiene
$1\,340$ hablantes y $6\,700$ clips. De los cinco clips
de cada hablante, tres se usan para la inscripción del ASV
y dos para las pruebas.

Los clips originales, en formato MP3, se convierten a
WAV y se remuestrean a $16\,\mathrm{kHz}$.

Common Voice no tiene un estrato colombiano aislable:
las voces colombianas quedan dentro de las categorías
«andino-pacífico» y «caribe». Esta es una limitación
del corpus para el contexto del proyecto.

\subsubsection{Métricas y criterios de evaluación}

Cada métrica responde a una parte de la pregunta de investigación. La Tabla~\ref{tab:metricas} resume qué mide cada una, cómo se obtiene, para qué se usa y cuál es el criterio del proyecto.

\begin{table*}[t]
\centering
\caption{Métricas de evaluación del proyecto.}
\label{tab:metricas}
\small
\renewcommand{\arraystretch}{1.2}
\begin{tabularx}{\textwidth}{
>{\raggedright\arraybackslash}p{2.6cm}
>{\raggedright\arraybackslash}p{3.0cm}
>{\raggedright\arraybackslash}X
>{\raggedright\arraybackslash}p{3.1cm}
>{\centering\arraybackslash}p{1.7cm}}
\toprule
\textbf{Métrica} & \textbf{Qué mide} & \textbf{Cómo se obtiene} & \textbf{Para qué se usa} & \textbf{Criterio} \\
\midrule
EER (privacidad)
& Qué tan bien reconoce un atacante al hablante después de anonimizar
& ASV ECAPA-TDNN preentrenado, en los escenarios o--o, o--a y a--a
& Determinar si la anonimización protege la identidad vocal
& EER $>30$\,\char37{} \\

WER (inteligibilidad)
& Cuánto contenido se pierde al anonimizar
& ASR Whisper \textit{turbo} preentrenado, sobre las grabaciones originales y anonimizadas
& Determinar si la conversación se sigue entendiendo
& WER $<15$\,\char37{} \\

RTF (tiempo real)
& Cuánto tarda el procesamiento frente a la duración del audio
& Tiempo de procesamiento del núcleo dividido por la duración del audio
& Determinar si el sistema puede operar durante una llamada
& $\mathrm{RTF}<0.1$ \\
\bottomrule
\end{tabularx}
\end{table*}

Los tres criterios son metas del proyecto. Una configuración cumple el objetivo del proyecto solo si satisface los tres a la vez.

\subsubsection{Diseño de Experimentos}

Los experimentos no buscan un único valor ideal del
coeficiente $\alpha$. Lo que interesa es entender cómo
responde el sistema cuando $\alpha$ cambia, qué intervalos
mantienen un buen equilibrio entre privacidad e
inteligibilidad y desde qué punto la voz empieza a volverse
difícil de entender. Es muy posible que varios intervalos
funcionen de manera parecida, y por eso el análisis apunta
a identificarlos, más que a quedarse con uno solo.

El trabajo experimental se hizo en tres etapas.

La primera fue un barrido con $\alpha$ fijo. Aquí todas
las grabaciones se anonimizaron con el mismo valor de
$\alpha$, probando valores de $0.2$ a $0.9$ en pasos de
$0.1$. También se incluyó $\alpha = 1.0$ como control,
porque con ese valor la transformación no cambia la voz
y los resultados deberían coincidir con los del audio
original. Esta etapa muestra el efecto de la intensidad
de la transformación por sí sola, sin la variación que
introduce el sorteo aleatorio, y permite ubicar los límites
del sistema: con $\alpha$ muy bajo la voz deja de entenderse,
y con $\alpha$ muy alto casi no hay protección. Como en este
caso las grabaciones de inscripción y de prueba quedan
transformadas exactamente igual, el escenario a--a no
refleja cómo funcionaría el sistema en la práctica, así
que el análisis de esta etapa se centra en el escenario o--a.

La segunda etapa fue un barrido más fino, también con
$\alpha$ fijo. Una vez ubicada la zona donde la privacidad
y la inteligibilidad cambian con más fuerza, se evaluaron
valores entre $0.65$ y $0.875$ en pasos de $0.025$, para
ver con más detalle qué ocurre en esa zona.

La tercera etapa evaluó intervalos uniformes de $\alpha$,
que es como funciona el sistema en su uso real: cada
grabación recibe un $\alpha$ elegido al azar dentro del
intervalo. Se probaron nueve intervalos, que resultan de
combinar tres puntos medios ($0.725$, $0.75$ y $0.775$)
con tres anchos ($0.18$, $0.15$ y $0.12$). En esta etapa
sí se analizan los dos escenarios, o--a y a--a.

El ancho del intervalo tiene una razón de ser. Para
deshacer la transformación y recuperar la voz original,
un atacante tendría que conocer el valor exacto de
$\alpha$ que se usó en cada grabación
\cite{patino2021}. Mientras más ancho es el intervalo,
más valores puede tomar $\alpha$ y mayor es su entropía,
es decir, más incierto resulta su valor y más difícil
es adivinarlo. Pero el intervalo no se puede ensanchar
indefinidamente, porque los valores bajos de $\alpha$
dañan la inteligibilidad y los altos protegen poco.
Por eso los anchos que se probaron están acotados por
los límites que mostró el barrido con $\alpha$ fijo.

En las tres etapas se miden la EER, como indicador de
privacidad, y la WER, como indicador de inteligibilidad.

Todas las configuraciones se evaluaron sobre las mismas
grabaciones. En el caso de los intervalos, el $\alpha$
de cada grabación se genera de forma determinista a
partir de su identificador, de modo que un mismo clip
recibe siempre el mismo sorteo. Gracias a esto, las
configuraciones se pueden comparar de forma pareada,
sobre exactamente los mismos datos.

La incertidumbre de cada resultado se estima con
intervalos de confianza calculados por \textit{bootstrap}.
Para comparar configuraciones vecinas se usa un
\textit{bootstrap} pareado de la diferencia entre ellas,
que indica si el cambio de una a otra es real o podría
explicarse por el azar del muestreo. Una parte del
análisis es también cualitativa: se observa cómo
evolucionan las métricas a medida que cambia $\alpha$
y se escuchan las grabaciones anonimizadas.

\subsection{Prototipo en tiempo real}

El prototipo reutiliza el mismo núcleo de procesamiento en
modo \textit{streaming}. Un adaptador entrega el audio al
núcleo en bloques de tamaño fijo, y se verificó que su salida
es equivalente a la del \textit{pipeline} por lotes.

El montaje usa la interfaz de audio MME de Windows a
$16\,\mathrm{kHz}$, con el búfer por defecto, un micrófono y
audífonos. El procedimiento consiste en tres ensayos de
$60\,\mathrm{s}$ de habla continua. En cada ensayo se mide:

\begin{itemize}
    \item El RTF, como el tiempo de procesamiento de cada
    bloque dividido por la duración del bloque, promediado
    sobre el ensayo.

    \item La latencia estimada, como la suma del retardo del
    búfer y el retardo algorítmico del procesamiento por
    tramas.

    \item Los cortes de audio durante el ensayo.
\end{itemize}


Como marco de referencia para evaluar la latencia se
utiliza la Recomendación UIT-T G.114
\cite{ituG114}, que establece orientaciones sobre
el retardo unidireccional en comunicaciones de voz.
Según esta recomendación, cuando el retardo total
de extremo a extremo es inferior a
$150\,\mathrm{ms}$, la mayoría de las aplicaciones
presentan una interactividad prácticamente
transparente. Asimismo, se establece un límite
superior recomendado de $400\,\mathrm{ms}$ para
la planificación general de redes.

En este trabajo, el valor de $150\,\mathrm{ms}$
se utiliza como referencia para interpretar 
la latencia estimada del prototipo.
Sin embargo, esta medición no representa el
retardo total de una llamada telefónica, puesto
que no incluye los retardos asociados con la
transmisión por la red celular.


\subsection{Consideraciones éticas}

El proyecto usa solo corpus públicos y no recolecta voces
de terceros; en las pruebas del prototipo se usa únicamente
la voz de la autora. El Protocolo de Datos clasificó el manejo
de los datos como ``Anónimo'', y el proyecto se clasificó
como ``Sin riesgo'' según el diagrama de autoclasificación
del Comité de Ética de la Facultad de Ingeniería.

Como decisión de diseño, el valor de $\alpha$ de cada sesión
no se registra junto al identificador de la sesión. Si ambos
se guardaran juntos, ya que, si alguien accediera al registro, podría
recuperar el coeficiente y facilitar la reversión de la
transformación.

% ================================================================
% 4. IMPLEMENTACION
% ================================================================
\section{Implementación}

% ================================================================
% 5. RESULTADOS
% ================================================================
\section{Resultados}

% ================================================================
% 6. CONCLUSIONES
% ================================================================
\section{Conclusiones}

% ================================================================
% TABLA DE ACRONIMOS
% ================================================================
\section*{Tabla de acrónimos}

\begin{table*}[t]
\centering
\small
\renewcommand{\arraystretch}{1.15}
\begin{tabularx}{\textwidth}{>{\raggedright\arraybackslash}p{3.2cm}X}
\toprule
\textbf{Acrónimo} & \textbf{Significado} \\
\midrule
a--a, o--a, o--o & Escenarios de evaluación: anonimizado--anonimizado, original--anonimizado, original--original \\
ASR & Reconocimiento automático de habla (\textit{Automatic Speech Recognition}) \\
ASV & Verificación automática de locutor (\textit{Automatic Speaker Verification}) \\
CPU & Unidad central de procesamiento (\textit{Central Processing Unit}) \\
ECAPA-TDNN & \textit{Emphasized Channel Attention, Propagation and Aggregation in Time Delay Neural Network} \\
EER & Tasa de error (\textit{Equal Error Rate}) \\
FAR & Tasa de falsa aceptación (\textit{False Acceptance Rate}) \\
FRR & Tasa de falso rechazo (\textit{False Rejection Rate}) \\
GPU & Unidad de procesamiento gráfico (\textit{Graphics Processing Unit}) \\
IA & Inteligencia artificial \\
LPC & Predicción lineal (\textit{Linear Predictive Coding}) \\
RTC & Coeficiente de tiempo real (\textit{Real-Time Coefficient}) \\
RTF & Factor de tiempo real (\textit{Real-Time Factor}) \\
VPC & \textit{VoicePrivacy Challenge} \\
WAV & \textit{Waveform Audio File Format} \\
WER & Tasa de error de palabras (\textit{Word Error Rate}) \\
ZEBRA & \textit{Zero Evidence Biometric Recognition Assessment} \\
\bottomrule
\end{tabularx}
\end{table*}

% ================================================================
% REFERENCIAS
% ================================================================
\begin{thebibliography}{99}

\bibitem{patino2021}
J. Patino, N. Tomashenko, M. Todisco, A. Nautsch, and N. Evans,
``Speaker anonymisation using the McAdams coefficient,''
\textit{Proceedings of Interspeech 2021}, pp. 1099--1103, 2021.
doi: 10.21437/Interspeech.2021-1070.

\bibitem{deng2023}
J. Deng, F. Teng, Y. Chen, X. Chen, Z. Wang, and W. Xu,
``V-Cloak: Intelligibility-, naturalness- \& timbre-preserving real-time voice anonymization,''
\textit{32nd USENIX Security Symposium}, pp. 5181--5198, 2023.
\url{https://arxiv.org/abs/2210.15140}

\bibitem{quamer2024}
W. Quamer and R. Gutierrez-Osuna,
``End-to-end streaming model for low-latency speech anonymization,''
\textit{2024 IEEE Spoken Language Technology Workshop (SLT)}, pp. 727--734, 2024.
\url{https://arxiv.org/abs/2406.09277}

\bibitem{tomashenko2024}
N. Tomashenko et al.,
``The VoicePrivacy 2024 Challenge,''
\textit{4th Symposium on Security and Privacy in Speech Communication}, Kos, Greece, 2024.
\url{https://www.voiceprivacychallenge.org}

\bibitem{kuzmin2026}
N. Kuzmin, S. Liu, K. A. Lee, and E. S. Chng,
``Stream-Voice-Anon: Enhancing utility of real-time speaker anonymization via neural audio codec and language models,''
\textit{arXiv preprint arXiv:2601.13948}, 2026.
\url{https://arxiv.org/abs/2601.13948}

\bibitem{markel1976linear}
J. D. Markel and A. H. Gray Jr.,
\textit{Linear Prediction of Speech}.
Berlin, Heidelberg: Springer, 1976.
doi: 10.1007/978-3-642-66286-7.

\bibitem{tayebi2024}
S. Tayebi Arasteh, T. Arias-Vergara, P. A. Pérez-Toro, T. Weise, K. Packhäuser, M. Schuster, E. Nöth, A. Maier, and S. H. Yang,
``Addressing challenges in speaker anonymization to maintain utility while ensuring privacy of pathological speech,''
\textit{Communications Medicine}, vol. 4, Art. 182, 2024.
doi: 10.1038/s43856-024-00609-5.

\bibitem{ardila2020}
R. Ardila, M. Branson, K. Davis, M. Kohler, J. Meyer, M. Henretty, R. Morais, L. Saunders, F. Tyers, and G. Weber,
``Common Voice: A massively-multilingual speech corpus,''
\textit{Proceedings of the Twelfth Language Resources and Evaluation Conference}, pp. 4218--4222, 2020.
\url{https://aclanthology.org/2020.lrec-1.520/}

\bibitem{acis2026}
ACIS,
``Colombianos recibieron más de 16.630 millones de llamadas spam durante 2025,'' 9 de febrero de 2026.
\url{https://www.acis.org.co/blog/noticias-2/colombianos-recibieron-mas-de-16-630-millones-de-llamadas-spam-durante-2025-6036}

\bibitem{barrington2025}
S. Barrington, E. A. Cooper, and H. Farid,
``People are poorly equipped to detect AI-powered voice clones,''
\textit{Scientific Reports}, vol. 15, Art. 11004, 2025.
\url{https://doi.org/10.1038/s41598-025-94170-3}

\bibitem{groupib2025}
Group-IB,
\textit{The Voice of Fraud: Deepfake Vishing and the New Age of Social Engineering}, informe técnico, 2025.
\url{https://www.group-ib.com/resources/research-hub/voice-of-fraud/}

\bibitem{transunion2025}
TransUnion Colombia,
``27\% de los intentos de fraude digital en Colombia son a través de llamadas telefónicas,'' 22 de abril de 2025.
\url{https://noticias.transunion.co/27-de-los-intentos-de-fraude-digital-en-colombia-son-a-traves-de-llamadas-telefonicas/}

\bibitem{desplanques2020}
B. Desplanques, J. Thienpondt, and K. Demuynck,
``ECAPA-TDNN: Emphasized channel attention, propagation and aggregation in TDNN based speaker verification,''
\textit{Proceedings of Interspeech 2020}, pp. 3830--3834, 2020.
doi: 10.21437/Interspeech.2020-2650.

\bibitem{radford2023}
A. Radford, J. W. Kim, T. Xu, G. Brockman, C. McLeavey, and I. Sutskever,
``Robust speech recognition via large-scale weak supervision,''
\textit{Proceedings of the 40th International Conference on Machine Learning}, PMLR, vol. 202, pp. 28492--28518, 2023.
\url{https://proceedings.mlr.press/v202/radford23a.html}


\bibitem{ituG114}
Unión Internacional de Telecomunicaciones,
\textit{Tiempo de transmisión en un sentido},
Recomendación UIT-T G.114 (05/2003).
Ginebra, Suiza: UIT, 2003.
[En línea]. Disponible en:
\url{https://www.itu.int/rec/T-REC-G.114/es}

\end{thebibliography}

\end{document}
